\section{The FSA-v2 SDE model}\label{sec:fsa}

\subsection{State, drift, and diffusion}

Let $(B(t),\,F(t),\,A(t))$ denote the latent state of a single individual
at time $t \ge 0$, measured in days. The three components are interpreted as
\begin{itemize}
  \item $B(t) \in [0,1]$ -- a normalised \emph{fitness} variable; the
        Banister chronic compartment;
  \item $F(t) \in [0,\infty)$ -- an unnormalised \emph{fatigue} variable;
        the Banister acute compartment;
  \item $A(t) \in [0,\infty)$ -- the \emph{amplitude} of an autonomic /
        circadian Stuart--Landau oscillator.
\end{itemize}
The single exogenous control input is the (non-negative) training-stimulus
rate $\Phi(t) \in [0,\Phimax]$, expressed in arbitrary daily TRIMP-like
units.

The model is presented in the \emph{G1-reparametrized} form. As we will
see in Section~\ref{sec:fim}, an earlier form of this model contained
three pairs of parameters that the four-channel observation model could
not separate. The G1 form rotates each pair into one strongly identified
parameter and one residual; the drift formulas below are
mathematically equivalent to the older form at a chosen operating point
$(B_{\mathrm{typ}},F_{\mathrm{typ}},A_{\mathrm{typ}})$ that we now fix.

\begin{definition}[Operating point]\label{def:operating-point}
We work throughout at the reference operating point
\begin{equation}
  A_{\mathrm{typ}} = 0.10, \qquad
  F_{\mathrm{typ}} = 0.20, \qquad
  \Phi_{\mathrm{typ}} = 1.0.
  \label{eq:operating-point}
\end{equation}
This is the regime of a de-trained athlete training at the canonical
Banister default load. The reparametrization is centred at this point,
and prior beliefs on the residual parameters encode ``the system is
typically near here.''
\end{definition}

The drift coefficients are
\begin{align}
  \mu(B,F) &= \mu_0 + \mu_B B - \mu_F F - \mu_{FF}\,(F - F_{\mathrm{typ}})^2,
            \label{eq:fsa-mu} \\
  \dd B &= \left[\,\kappa_B \cdot \frac{1+\varepsilon_A A}{1+\varepsilon_A A_{\mathrm{typ}}}\,\Phi(t)
                  - \frac{B}{\tau_B}\,\right]\dd t
            + \sigma_B \sqrt{B(1-B)}\,\dd W_B, \label{eq:fsa-B} \\
  \dd F &= \left[\,\kappa_F\,\Phi(t)
                  - \frac{1+\lambda_A A}{1+\lambda_A A_{\mathrm{typ}}}\,
                    \frac{F}{\tau_F}\,\right]\dd t
            + \sigma_F \sqrt{F}\,\dd W_F, \label{eq:fsa-F} \\
  \dd A &= \bigl[\,\mu(B,F)\,A - \eta\,A^3\,\bigr]\,\dd t
            + \sigma_A \sqrt{A}\,\dd W_A, \label{eq:fsa-A}
\end{align}
where $W_B,W_F,W_A$ are independent scalar Brownian motions.

The diffusions in \eqref{eq:fsa-B}--\eqref{eq:fsa-F} are unchanged from
the pre-G1 model: a Jacobi diffusion on $[0,1]$ for $B$ (Forman and
S\o rensen, 2008) whose volatility vanishes at both endpoints, and
square-root (CIR-type, Cox--Ingersoll--Ross 1985) diffusions for $F$ and
$A$ whose volatility vanishes at the lower boundary.

\paragraph{Reading the residual factors.} The two ratios
$(1+\varepsilon_A A)/(1+\varepsilon_A A_{\mathrm{typ}})$ and
$(1+\lambda_A A)/(1+\lambda_A A_{\mathrm{typ}})$ each equal one when
$A=A_{\mathrm{typ}}$. So at the operating point, and for the constant
control $\Phi \equiv \Phi_{\mathrm{typ}}$, the drift collapses to
\begin{equation}
  \dot B = \kappa_B \Phi - B/\tau_B,
  \qquad
  \dot F = \kappa_F \Phi - F/\tau_F,
  \qquad
  \mu = \mu_0 + \mu_B B - \mu_F F.
  \label{eq:fsa-collapse}
\end{equation}
That is, the strongly identified parameters $\kappa_B, \tau_F, \mu_F$
play exactly the role they would in a Banister model linearised at
$(F_{\mathrm{typ}}, A_{\mathrm{typ}})$. The residuals
$\varepsilon_A, \lambda_A, \mu_{FF}$ control departures from the
operating point, and they do so via terms that vanish at the operating
point itself; this is what makes them weakly identified at one-day
windows and consequently candidates for tighter priors.

\subsection{Truth parameters}

Table~\ref{tab:fsa-truth} lists the truth parameters of the
G1-reparametrized model. Compared to the pre-G1 form, three values
change: $\kappa_B = 0.012 \cdot (1+\varepsilon_A A_{\mathrm{typ}}) = 0.01248$,
$\tau_F = 7.0 / (1+\lambda_A A_{\mathrm{typ}}) = 6.\overline{36}\,$d,
and $\mu_F = 0.10 + 2 F_{\mathrm{typ}} \mu_{FF} = 0.26$, with a small
additive shift $\mu_0 \mapsto \mu_0 + \mu_{FF} F_{\mathrm{typ}}^2 = 0.036$
absorbing the constant from completing the square in
\eqref{eq:fsa-mu}.

\begin{table}[h]
\centering
\small
\begin{tabular}{@{}llrl@{}}
\toprule
Group & Symbol & Value & Meaning \\
\midrule
Banister timescales       & $\tau_B$              & $42.000$    & chronic time constant (d) \\
                          & $\tau_F$              & $6.364$     & effective acute time constant (d) \\
Banister gains            & $\kappa_B$            & $0.01248$   & effective $B$-gain at $A_{\mathrm{typ}}$ \\
                          & $\kappa_F$            & $0.030$     & $F$-gain per unit $\Phi$ \\
$A$-coupling (residual)   & $\varepsilon_A$       & $0.40$      & residual $B$-gain boost beyond $A_{\mathrm{typ}}$ \\
                          & $\lambda_A$           & $1.00$      & residual $F$-clearance boost beyond $A_{\mathrm{typ}}$ \\
Stuart--Landau            & $\mu_0$               & $0.036$     & growth-rate intercept \\
                          & $\mu_B$               & $0.30$      & fitness raises $\mu$ \\
                          & $\mu_F$               & $0.26$      & local slope of $\mu$ at $F_{\mathrm{typ}}$ \\
                          & $\mu_{FF}$            & $0.40$      & residual curvature of $\mu$ in $F$ \\
                          & $\eta$                & $0.20$      & cubic damping \\
Diffusion                 & $\sigma_B$            & $0.010$     & Jacobi noise scale \\
                          & $\sigma_F$            & $0.012$     & CIR noise scale on $F$ \\
                          & $\sigma_A$            & $0.020$     & CIR noise scale on $A$ \\
\bottomrule
\end{tabular}
\caption{Truth parameters of the G1-reparametrized FSA-v2 dynamics.
Source: \texttt{version\_2/models/fsa\_high\_res/\_dynamics.py}.}
\label{tab:fsa-truth}
\end{table}

\subsection{The training-stimulus envelope}

The control variable in the SDE is the instantaneous training-stimulus
rate $\Phi(t)$. In practice training is delivered in discrete daily
sessions concentrated in the morning, and the simulator expands a daily
schedule $(\Phi_d)_{d=1,\dots,N}$ into a high-resolution envelope at
fifteen-minute resolution as follows. Let
$\tau_{\mathrm{burst}}=3\,\mathrm{h}$ be a shape parameter, and let
$[t_{\mathrm{wake}},t_{\mathrm{sleep}}] = [7,23]$ be the daily wake
window (in hours). The within-day envelope is
\begin{equation}
  e(t) \;\propto\;
    \begin{cases}
      (t-t_{\mathrm{wake}})\,\exp\!\bigl(-(t-t_{\mathrm{wake}})/\tau_{\mathrm{burst}}\bigr)
        & t_{\mathrm{wake}} \le t < t_{\mathrm{sleep}}, \\
      0  & \text{otherwise},
    \end{cases}
  \label{eq:phi-envelope}
\end{equation}
which is the standard Gamma$(k=2)$ shape, peaking at
$t-t_{\mathrm{wake}} = \tau_{\mathrm{burst}}$. The constant of
proportionality is fixed so that
$\int_{\text{day}} \Phi(t)\,\dd t = 24\,\Phi_d$, i.e.~the daily integral
of the within-day envelope is preserved. The slow Banister equations
\eqref{eq:fsa-B}--\eqref{eq:fsa-F} therefore see the same daily load
whether $\Phi$ is delivered as a constant or as a morning-loaded burst,
and the fast Stuart--Landau dynamics in $A$ pick up the within-day shape.

\subsection{Four-channel observation model}

The state $(B,F,A)$ is not directly observed. Instead we receive four
discrete-time channels at the fifteen-minute grid:
\begin{align}
  \mathrm{HR}(t) &\sim
    \Ncal\!\bigl(\,\mathrm{HR}_{\mathrm{base}}
                  - \kappa_B^{\mathrm{HR}}\,B
                  + \alpha_A^{\mathrm{HR}}\,A
                  + \beta_C^{\mathrm{HR}}\,C(t),\;
                  \sigma_{\mathrm{HR}}^2\bigr),
    && \text{(sleep-gated)} \label{eq:obs-hr}\\
  P(\mathrm{sleep}(t)=1\mid A,C) &=
    \sigmoid\!\bigl(\,k_C\,C(t) + k_A\,A - \tilde c\,\bigr),
    && \text{(always observed)} \label{eq:obs-sleep}\\
  S(t) &\sim
    \Ncal\!\bigl(\,S_{\mathrm{base}} + k_F\,F - k_A^S\,A
                  + \beta_C^S\,C(t),\;
                  \sigma_S^2\bigr),
    && \text{(wake-gated)} \label{eq:obs-stress}\\
  \log(\mathrm{steps}(t)+1) &\sim
    \Ncal\!\bigl(\,\mu_0^{\mathrm{step}} + \beta_B^{\mathrm{step}}\,B
                  - \beta_F^{\mathrm{step}}\,F + \beta_A^{\mathrm{step}}\,A
                  + \beta_C^{\mathrm{step}}\,C(t),\;
                  \sigma_{\mathrm{step}}^2\bigr),
    && \text{(wake-gated)} \label{eq:obs-steps}
\end{align}
where $C(t) = \cos(2\pi t + \phi)$ is a fixed deterministic circadian
covariate (not a latent variable). ``Sleep-gated'' channels are emitted
only when the sleep label of \eqref{eq:obs-sleep} is one; ``wake-gated''
channels only when it is zero. The collection of observation-model
parameters $(\mathrm{HR}_{\mathrm{base}},\kappa_B^{\mathrm{HR}},\dots,\sigma_{\mathrm{step}})$
brings the full static parameter vector $\thetadyn$ for the joint
state-space model to $30$ scalar entries. We write $y_{1:n}$ for the
multi-channel observation record up to time $t_n$ and
$p(y_{1:n}\mid \thetadyn)$ for its marginal likelihood under the model.

\medskip
\noindent
The two analyses that an earlier version of these notes flagged as
outstanding -- Fisher Information rank identifiability and Lyapunov
stability -- have since been carried out and form
Sections~\ref{sec:fim} and~\ref{sec:lyap} respectively.
